% !TeX root = ../../Book4.tex
% 纲要标题：### 17.3 上同调函子
% 纲要标签：b4:sec:1603
% 纲要位置：计划纲要.md，第 3668 行。

\section{上同调函子}
\label{b4:sec:1603}

好三角经过上同调函子以后形成长正合列，Hom 本身也具有这种性质。这些正合列判断三角何时分裂，并说明它与交换范畴短正合列的关系。


% 纲要内容（保留纲要核验项目）：
%
% * 三角产生长正合列
% * Hom 函子的正合性
% * 三角与短正合列的对应范围
%

\subsection{三角与长正合列}
\label{b4:subsec:1603:01}

\begin{definition}\label{b4:def:cohomological-functor}
设 \(\mathcal T\) 为三角范畴，\(\mathcal B\) 为交换范畴。加性函子 \(H:\mathcal T\to\mathcal B\) 若将每个好三角 \(X\to Y\to Z\to X[1]\) 变为在中间正合的序列 \(H(X)\to H(Y)\to H(Z)\)，则称为\term[shangtongdiao-hanzi]{上同调函子}[cohomological functor]。记 \(H^n(X)=H(X[n])\)。
\end{definition}

由旋转和平移，定义中的三项正合自动延长成
\[
 \cdots\to H^n(X)\to H^n(Y)\to H^n(Z)
 \to H^{n+1}(X)\to H^{n+1}(Y)\to\cdots.
\]
旋转可能给某些箭头添负号，但不改变核和像；选定一致的符号后，这就是长正合列。三角态射诱导长正合列之间的交换图，所以正合列同时记录对象和态射的信息。

\begin{proposition}\label{b4:prop:cohomology-on-homotopy}
设 \(\mathcal A\) 为局部小交换范畴。通常的零次上同调 \(H^0:K(\mathcal A)\to\mathcal A\) 是上同调函子，其平移得到通常的各次上同调。
\end{proposition}
\begin{proof}
链同伦映射诱导相同上同调，故函子已在 \(K(\mathcal A)\) 上定义。对标准锥三角，正合列由逐次分裂短正合列 \(0\to B\to C_f\to A[1]\to0\) 得到，连接映射是 \(H^{n+1}(f)\)，见\cref{b4:prop:cone-long-exact}。因此其中的 \(H^n(A)\to H^n(B)\to H^n(C_f)\) 正合。好三角与标准锥三角同构，正合性随同构保留。
\end{proof}

\subsection{Hom 函子的正合性}
\label{b4:subsec:1603:02}

\begin{proposition}\label{b4:prop:triangle-hom-exact}
设 \(\mathcal T\) 为局部小三角范畴，\(W\in\mathcal T\)。函子 \(\operatorname{Hom}_{\mathcal T}(W,-)\) 是取值于交换群的上同调函子；\(\operatorname{Hom}_{\mathcal T}(-,W)\) 具有相应的反变正合性。特别地，好三角 \(X\to Y\to Z\to X[1]\) 给出
\[
 \cdots\to\operatorname{Hom}(W,X[n])\to\operatorname{Hom}(W,Y[n])
 \to\operatorname{Hom}(W,Z[n])\to\operatorname{Hom}(W,X[n+1])\to\cdots,
\]
以及箭头倒向的反变长正合列。
\end{proposition}
\begin{proof}
\cref{b4:lem:pretriangle-hom}已证明三项正合。对三角旋转并平移，可以把长列中的任意一项移到三项列的中间，故在每一项正合。反变情形应用该引理的第二列。
\end{proof}

\begin{proposition}\label{b4:prop:split-triangle-criterion}
设 \(\mathcal T\) 为局部小三角范畴，\(X\xrightarrow{u}Y\xrightarrow{v}Z\xrightarrow{w}X[1]\) 为好三角。则 \(w=0\) 当且仅当此三角同构于由包含、投影和零映射组成的分裂三角 \(X\to X\oplus Z\to Z\to X[1]\)。
\end{proposition}
\begin{proof}
若 \(w=0\)，协变 Hom 长正合列的片段
\[
 \operatorname{Hom}(Z,Y)\xrightarrow{v_*}\operatorname{Hom}(Z,Z)
 \xrightarrow{w_*}\operatorname{Hom}(Z,X[1])
\]
中最后一箭头为零，所以 \(1_Z\) 可提升为 \(s:Z\to Y\)，满足 \(vs=1_Z\)。另一方面，逆向旋转后取反变 Hom，得到
\[
 \operatorname{Hom}(Y,X)\xrightarrow{u^*}\operatorname{Hom}(X,X)
 \xrightarrow{(-w[-1])^*}\operatorname{Hom}(Z[-1],X).
\]
末箭头也为零，所以 \(1_X\) 可提升为 \(r:Y\to X\)，满足 \(ru=1_X\)。连接箭头消失因而同时提供了 \(u\) 的左逆和 \(v\) 的右逆。

要用它们构造 \(Y\cong X\oplus Z\)，还需要使两个分量相容，即 \(rs=0\)。以 \(s-urs\) 代替 \(s\)，由于 \(vu=0\) 且 \(ru=1_X\)，仍有 \(vs=1_Z\)，并且 \(rs\) 变为零。此时 \((r,v)(u,s)=1_{X\oplus Z}\)。为证明另一方向的复合也是恒等，令 \(e=1_Y-ur-sv\)，则 \(ve=0\)、\(re=0\)。对 \(\operatorname{Hom}(Y,-)\) 使用正合性，存在 \(t:Y\to X\) 使 \(e=ut\)；再由 \(t=rut=re=0\)，得 \(e=0\)。所以 \((u,s):X\oplus Z\to Y\) 与 \((r,v)\) 互逆，并保持三角的全部箭头。反向直接从分裂三角的最后箭头为零得到。
\end{proof}

这一判据说明，三角的连接箭头为零时，中间对象连同两端映射一起分裂；这与\cref{b4:thm:ext1-extensions}中短正合扩张的零类对应分裂扩张相呼应。

\ProseParagraphBreak
设 \(\mathcal A\) 为局部小加性范畴，\(C,D\) 为其中的复形。在 \(K(\mathcal A)\) 中还有一个计算形式：
\[
 \operatorname{Hom}_{K(\mathcal A)}(C,D[n])
 \cong H^n\operatorname{Hom}_{\mathcal A}^{\bullet}(C,D).
\]
右端按\cref{b4:def:hom-complex}的公式，把 Hom 换为 \(\mathcal A\) 的态射交换群。次数 \(n\) 的余圈条件为 \(\mathrm{d}_Df=(-1)^nf\mathrm{d}_C\)，恰是 \(C\to D[n]\) 的复形映射条件。还须比较两边所商掉的子群。若 \(g\) 是次数 \(n-1\) 的 Hom 元素，则
\[
 \mathrm{d}_{\operatorname{Hom}}g
 =\mathrm{d}_Dg-(-1)^{n-1}g\mathrm{d}_C
 =\mathrm{d}_Dg+(-1)^ng\mathrm{d}_C.
\]
把 \(h=(-1)^ng\) 看作从 \(C\) 到 \(D[n]\) 的降次态射族，得到
\[
\begin{aligned}
 \mathrm{d}_{D[n]}h+h\mathrm{d}_C
 &=(-1)^n\cdot\mathrm{d}_Dh+h\mathrm{d}_C\\
 &=\mathrm{d}_Dg+(-1)^ng\mathrm{d}_C
 =\mathrm{d}_{\operatorname{Hom}}g.
\end{aligned}
\]
反过来，每个这样的 \(h\) 都唯一写成 \((-1)^ng\)，故 Hom 复形的余边与到 \(D[n]\) 的零伦映射组成同一子群，取商便得到上述同构。此处必须取整个 Hom 复形的上同调，而不能只看某一个次数的普通 Hom 群。

\subsection{三角与短正合列的对应范围}
\label{b4:subsec:1603:03}

\begin{proposition}\label{b4:prop:degreewise-split-triangle}
设 \(\mathcal A\) 为局部小交换范畴，\(0\to A\xrightarrow{j}B\xrightarrow{q}C\to0\) 是逐次分裂的复形短正合列。则存在 \(\delta:C\to A[1]\)，使 \(A\xrightarrow{j}B\xrightarrow{q}C\xrightarrow{\delta}A[1]\) 是 \(K(\mathcal A)\) 中的好三角。
\end{proposition}
\begin{proof}
选逐次分裂，写 \(B^n=A^n\oplus C^n\)，微分为
\(\mathrm{d}_B=\left(\begin{smallmatrix}\mathrm{d}_A&t\\0&\mathrm{d}_C\end{smallmatrix}\right)\)。条件 \(\mathrm{d}_B^2=0\) 给出 \(\mathrm{d}_At+t\mathrm{d}_C=0\)，故 \(-t:C\to A[1]\) 为复形映射。自然映射
\[
 r:\operatorname{Cone}(j)\to C,\quad r((a,c),a')=c,
 \qquad s(c)=((0,c),-tc)
\]
满足 \(rs=1\)。取 \(h((a,c),a')=((0,0),a)\)，便有 \(1-sr=\mathrm{d}h+h\mathrm{d}\)。于是 \(r\) 为同伦等价，其与锥的包含复合是 \(q\)。记锥的最后投影为 \(p_j:\operatorname{Cone}(j)\to A[1]\)，则 \(p_js=-t\)。由于 \(sr=1\) 在同伦范畴中成立，还有
\[
 (-t)r=p_jsr=p_j\qquad\text{于}\;K(\mathcal A).
\]
这同时核对了三角同构的最后一个方块。因此把锥三角经 \(r\) 运送，得到 \(\delta=-t\)。
\end{proof}

为比较这里的 \(\delta=-t\) 与普通短正合列的连接同态，记 \(z_C^n:Z^n(C)\to C^n\) 为余圈嵌入。逐次截面将它送到 \(B^n\) 的第二分量；再作微分得到 \((t^nz_C^n,0)\)。由\cref{b4:prop:connecting-map-construction}，\(t^nz_C^n\) 在上同调上诱导通常的连接同态。在模中，这就是把余圈 \(c\) 提升为 \((0,c)\) 并计算
\[
 \mathrm{d}_B(0,c)=(t(c),0),\qquad
 \delta_{\mathrm{LES}}[c]=[t(c)],\qquad
 H^n(\delta)[c]=-[t(c)].
\]
这里使用 \(H^n(A[1])=H^{n+1}(A)\) 的通常识别。因此三角连接箭头在上同调上的作用为 \(-\delta_{\mathrm{LES}}\)；负号来自本书把锥三角最后箭头固定为正投影的约定。它不影响正合性，却必须在比较连接同态的具体公式时保留。

\ProseParagraphBreak
若短正合列不逐次分裂，则不能直接使用这个结论。例如次数零上的 \(0\to\mathbb Z\xrightarrow{2}\mathbb Z\to\mathbb Z/2\mathbb Z\to0\)，\(\operatorname{Cone}(2)\to\mathbb Z/2\mathbb Z\) 是拟同构，却不是同伦等价：其同伦逆会给出从有限群到自由交换群的非零映射。到导出范畴中，这个拟同构才成为同构，短正合列也才无须逐次分裂即可给出好三角，见\cref{b4:prop:derived-short-exact-triangle}。
